\documentclass[11pt]{article}
\usepackage{amsmath,amssymb,amsthm,mathtools}
\usepackage{geometry}
\geometry{margin=1in}
\newtheorem{theorem}{Theorem}
\newtheorem{proposition}{Proposition}
\newtheorem{lemma}{Lemma}
\newtheorem{definition}{Definition}
\newtheorem{remark}{Remark}
\newtheorem{corollary}{Corollary}

\title{Step 54: Navier Strict-Extension Adequacy Mechanisms and Closeout}
\author{Anti-Localization / Formed-Layer Membrane Theory}
\date{}

\begin{document}
\maketitle

\section{Purpose and scope}
This note closes the Navier--Stokes aside at the level appropriate for the formed-layer membrane paper. It is not a Navier--Stokes regularity proof and it is not a Six Birds simulation. Its purpose is to classify, in the language of the adequacy residual
\[
\Xi_C(D\mid L),
\]
which lawful strict extensions could reduce the blind spot between native energy/enstrophy ledgers and layer-dissolving blow-up probes.

The conclusion is that classical energy and enstrophy are not adequate for pointwise blow-up probes in three dimensions. A Navier membrane proof must therefore supply one of several explicit strict-extension records: higher-order coercive audit, channelized feasible carrier, parabolic smoothing with source-admission, geometric/nonlinear source constraint, or an honest scoped nonclaim.

\section{Navier membrane setup}
At a scale or shell level $j$, let
\[
E_j
\]
be the exact packaged flow carrier, for example a divergence-free shell carrier or a finite Galerkin/shell package after all constraints are declared. Let
\[
C_j\succeq0
\]
be the audit energy, such as energy, enstrophy, dissipation, or a stronger Sobolev audit. Let
\[
L_j:E_j\to Y_j
\]
be the native probe family supplied by the formed flow layer, and let
\[
D_j:E_j\to Z_j
\]
be the layer-dissolving probe family, such as pointwise velocity, vorticity, strain, gradient, or local blow-up-rate probes.

On the legal quotient $(\ker C_j)^\perp$, define
\[
\Xi_{C_j}(D_j\mid L_j)
=
K_{DD,j}-K_{DL,j}K_{LL,j}^{\dagger}K_{LD,j}.
\]
This is the blind-spot currency of the blow-up probes after conditioning on the native flow probes.

\section{Sobolev scaling obstruction}
On $\mathbb T^d$ with Fourier cutoff $N$, let
\[
C_{s,N}e_k=(1+|k|^2)^s e_k
\]
and let $D_{q,x,N}$ be a pointwise derivative probe of order $q$. Then
\[
\operatorname{Cap}_{s,N}(D_q)
=\sum_{|k|\le N}\frac{|k|^{2q}}{(1+|k|^2)^s}.
\]

\begin{proposition}[Sobolev blow-up probe threshold]
As $N\to\infty$,
\[
\operatorname{Cap}_{s,N}(D_q)<\infty \text{ uniformly in }N
\]
if and only if
\[
s>q+\frac d2.
\]
At the endpoint the capacity grows logarithmically, and below the endpoint it grows like
\[
N^{d+2q-2s}.
\]
\end{proposition}

\begin{corollary}[Three-dimensional obstruction]
For $d=3$, pointwise velocity probes require $s>3/2$, and pointwise gradient/vorticity/strain probes require
\[
s>5/2.
\]
Thus the standard energy audit $s=0$ and enstrophy/dissipation audit $s=1$ are not adequate for pointwise blow-up probes.
\end{corollary}

\section{Strict-extension mechanisms}
\subsection{Native-family extension}
If a strict extension adds native probes
\[
M_j:E_j\to W_j,
\qquad
L_j^+=\begin{bmatrix}L_j\\ M_j\end{bmatrix},
\]
then
\[
\Xi_{C_j}(D_j\mid L_j^+)
=
\Xi_{C_j}(D_j\mid L_j)
-
K_{DM\mid L,j}K_{MM\mid L,j}^{\dagger}K_{MD\mid L,j}.
\]
Hence
\[
\Xi_{C_j}(D_j\mid L_j^+)\preceq \Xi_{C_j}(D_j\mid L_j).
\]
The blind spot decreases strictly only when the new native probes see the old residual.

\subsection{Higher-order coercive audit}
A stronger audit can bound the blow-up probes if it supplies the Sobolev threshold $s>q+d/2$ or an equivalent source-frame inequality
\[
C_j^+\succeq \lambda_jD_j^*\Theta_D^{-1}D_j.
\]
Then
\[
D_j(C_j^+)^{\dagger}D_j^*\preceq \lambda_j^{-1}\Theta_D.
\]
This is proof-strengthening only if the stronger audit is lawfully available from the flow closure. Adding an artificial high-order audit is not a Navier proof.

\subsection{Channelized feasible carrier}
If the exact flow carrier restricts feasible states to a channelized subcarrier
\[
\mathcal F_j=\operatorname{Ran}\Gamma_j
\]
with lower frame/energy bounds and delocalized channels, then pointwise probes may be controlled even when the ambient Sobolev audit is too weak. This is a genuine strict-extension mechanism only if the channelization is produced by the lawful dynamics or closure, not imposed post hoc.

\subsection{Parabolic smoothing bridge}
For the heat semigroup,
\[
\|\nabla^q e^{\tau\nu\Delta}u\|_{L^\infty}
\lesssim
\tau^{-\frac{q+d/2-s}{2}}\|u\|_{H^s}.
\]
This supplies a lagged smoothing bridge for $\tau>0$, but it does not by itself control instantaneous blow-up as $\tau\downarrow0$. It becomes a membrane mechanism only with source-admission and predictive transport records preventing nonlinear refocusing from reopening the blind spot.

\subsection{Nonlinear source-admission and geometric constraints}
Vorticity alignment, depletion of vortex stretching, pressure cancellation, or scale-local flux restrictions may reduce $\Xi$ only if they are expressed as an adequacy residual bound
\[
D_j=A_jL_j+R_j,
\qquad
R_jC_j^{\dagger}R_j^*\preceq \Xi_j,
\]
with predictive control of $\Xi_j$. Without this residual inequality, they are support evidence, not membrane proof.

\subsection{Scoped nonclaim / gating}
Restricting to a class already known to have stronger regularity, such as a high-Sobolev or Besov class, can produce a valid scoped membrane. It does not prove full Navier regularity unless the dynamics lawfully keeps all solutions in that class.

\section{Navier strict-extension adequacy theorem}
\begin{theorem}[Navier membrane obligation theorem]
Let $(E_j,C_j,L_j,D_j,\Gamma_j,R_j)$ be a formed predictive flow-layer record. Suppose:
\begin{enumerate}
\item the native membrane is accepted:
\[
K^L_j\preceq\Theta^L_j;
\]
\item the blow-up adequacy residual is controlled:
\[
\Xi_{C_j}(D_j\mid L_j)\preceq\Xi^D_j;
\]
\item the induced dissolving budget satisfies, after response transport,
\[
R_j\left(A_{*,j}\Theta^L_jA_{*,j}^*+\Xi^D_j\right)R_j^*\preceq\Theta^D
\]
for a common predictive budget $\Theta^D$;
\item all-six channel records are accepted, including null-mode legality, route-union honesty, packaging/canonicalization, predictive transport, and audit/currency records.
\end{enumerate}
Then no declared layer-dissolving Navier probe has a predictive native needle:
\[
\sup_j\operatorname{Cap}_{C_j}(y^*R_jD_j;E_j)\le y^*\Theta^D y
\qquad(y\in Z).
\]
\end{theorem}

\begin{remark}[Nonclaim]
The theorem does not prove Navier regularity. It states the exact membrane obligation a Navier proof must earn. Classical energy/enstrophy estimates alone fail the adequacy requirement for pointwise blow-up probes in dimension three.
\end{remark}

\section{Closeout for the present paper}
For the formed-layer membrane manuscript, Navier should be treated as a sanity-check application and obligation map. The main paper should not pursue a Navier proof. The load-bearing general object is the adequacy residual $\Xi_C(D\mid L)$, not any specific PDE calculation.

\end{document}
